% !TeX root = ../papers/nc-p3.tex
\section{半単純な判別法線を指定した族の最小周期}
\label{sec:marked-semisimple-period}
\begin{theorem}\label{thm:marked-semisimple-period}
定理\ref{thm:marked-prescribed-discriminant-as}の半単純な法線 $N$ と任意の直線 $\ell$ に対し、
そこで構成した既約族の幾何学的一般点の最小周期は4である。
印を忘れた場合も各位置への射影作用を取った添字保持自己同型群は自明である。
両因子を動かす十二次元の同型類の族の一般点でも成立する。
\end{theorem}
\begin{proof}
$N$ を対角化し、前節の $G_i=D_i$, $F_i=-D_i$, $Q_i=-[e_i]_\times$ を基点にする。
$B$ を位置1の配置、$A,C$ を位置0,2の配置とする。$b=(r,v)\in B$ の両隣
$a=(s,u)\in A$, $c=(h,w)\in C$ はそれぞれ
\[
 (sI+G_u)v=(rI+Q_v^t)u=0,\qquad
 (hI+F_w)v=(rI+Q_v)w=0
\]
で求まる。基点の $B$ は星点、$r=1$ で三成分のうち一つが零で他の二つが
$\pm\sqrt{-1/2}$ の十二点、$r=0$ で三座標点と $(1,\pm1,\pm1)$ の四点である。
四辺は二十個の被約点のグラフなので、基点の厳密局所化の近くでこれらを別々に標識できる。

一般点の射影同値は、標識のある近傍で、ある一つの置換に対して恒等的に成り立つ。
これを基点で調べる際の行列の退化を除く。非零行列の射影空間内で同値のグラフを閉包し、
各標識点について像と相手の点の二次小行列式を零とする。
基点へ特殊化した行列が特異なら、核に入らない点は像内に写る。
基点のどの平面も高々六点しか含まないので、核と像を合わせても二十点を対応させられない。
よって基点でも可逆である。基点の四つの座標点と全成分非零の一点は射影枠であり、
標識置換を固定すれば近傍で行列はこの五点から正則に復元できる。従って同じ置換が
\emph{全ての}接方向に沿って一次まで持ち上がらなければならない。

以下の小行列計算の入力と全検査は
\path{computations/verify_semisimple_period_obstruction.py}に置く。
分裂素数 $1009$ で還元し、二十点ずつの $A,C$ を上の二組の六式の核で構成する。
両者とも三点が共線になることはなく、六点を超える平面もない。四つの座標点は
四点平面十二枚・六点平面六枚に属し、残りはそれぞれ三枚・九枚に属する。
従って体を拡大しても還元配置間の射影同値は単項行列である。
座標置換と全成分非零の一点の像で列挙すると $A\to A$ と $A\to C$ は各96個である。
この列挙は標数零の候補もすべて検査する。実際、上の核・像の個数による議論は
この素数での特殊化にも適用できる。二十点は還元で分離したままであり、
置換を固定した行列は両配置の枠から単元を分母に持つ式で復元できる。

まず $T$ の行を $(3,2,-1),(1,-3,4),(-1,2,4)$ とし、
$\dot G_i=\sum_jT_{ij}D_j$, $\dot F=-\dot G$ とする。
$G'$ を固定するこの方向は、全ての指定 $N,\ell$ の族で取れる。
$\dot Q$ は前節の可逆微分で復元し、各点の微分は上の六式を二つのアフィンチャートで
微分した六次連立一次方程式で求める。左右から得た $\dot b$ は一致する。
候補 $M$ と対応点 $Mx=\lambda y$ について必要な式は
\[
              (\dot Mx+M\dot x-\lambda\dot y)\bmod ky=0.
\]
行列 $\dot M$ の十六成分への六十式の係数階数は15である。右辺を加えても
階数15となるのは、自己同型では恒等、$A\to C$ では $J=\operatorname{diag}(1,-1,-1,-1)$ だけである。

次に $G$ を固定する。$\dot F$ の対角跡零部分は三列各二成分である。任意の六成分に対して
$\sum[D_i,\dot F_i]=0$ となるように $D_j$ の係数を補える。
$[D_2,D_3],[D_3,D_1],[D_1,D_2]$ が基底であり、補い方の差は
$\dot F_i=\sum K_{ij}D_j$ の対称行列 $K$ である。この差は $J$ の持上げを妨げない。
実際 $F=(-I+tK)G$ の係数行列は非退化対称であり、定理\ref{thm:marked-point-commutators}
の $U$ 側の合同変換により基点で恒等な基底変更で $F=-G$ に戻せる。
その基底で $Q$ は交代解の一意なエタール持上げであり、二組の六式は
$(r,v)\mapsto(r,-v)$ と $(s,u)\mapsto(s,-u)$ で交換される。
従って対称な補正の六列は標数零でも $\dot M$ の像内にある。
これに対し六個の対角成分の単位方向を上の交換子基底で補い、
$J$ の一次方程式の右辺に並べると、還元で拡大階数は21である。
係数側には常にスカラー核があるため、標数零でも余核内の六列は独立である。

指定した $N,\ell$ の族では $G'$ を固定して $H\subset F'\subset N^\perp$ の
射影直線上で $F'=G'$ から離れる接方向がある。前節の階数六の基底束によりこれは
実際に持ち上がる。$N^\perp/G'$ は非零の対角跡零方向なので、この接方向の六成分は非零である。
よって $J$ は持ち上がらず、一般の二つの偶配置は射影同値でない。
二位置移動にはこの同値が必要であり、標準四周期性から最小周期が従う。
位置0で恒等な射影自己同型はグラフに沿って各位置の全点を固定する。
基点の偶配置は射影枠を含む。奇配置の七点を固定すれば無限遠平面上の作用はスカラーであり、
星点と有限の非零点も固定するので全体がスカラーである。この条件は近傍で開である。
反対の奇配置も混合点で符号変更により同じ配置なので、全位置の射影作用が恒等となる。
\end{proof}
